% !PRINTED 145
\weberpage{145}
If we assume, as we are free to do, that $b \geqq a$, then the
expression on the right side of this inequality is greater
than $n$, and equal to $n$ only when $b = n - 1, a = 1$. In
this special case $j = n$ can occur, but only when
the $n - 1$ elements are permuted by $Q$ in all possible ways,
which is expressed in Theorem I.

The expressions (3) for the lower limit of $j$ are nothing
other than the binomial coefficients $B_a^{(n)}$, whose formation
shows at once that up to the middle, i.e. as long as $2a \leqq n$,
they form an \gesperrt{increasing sequence of numbers}.

For further use we shall draw from this the
conclusion:

$\alpha$) \gesperrt{If $a = 1$, that is, if the group $Q$ leaves one figure
unchanged, then its degree is a divisor of
$\Pi(n-1)$; but if $a > 1$, then the degree of $Q$ is
$\leqq \Pi(n-2)\,\Pi(2)$.}

The theorem which we further wish to prove now reads:

\gesperrt{II. Besides the alternating group there is no
transitive and primitive divisor $Q$ of $P$
whose index is $\leqq n$, except in the
two cases $n = 4, n = 6$.}

In the proof we make use of the theorems (\S.\,153,
9., 10. of the first volume) that a transitive and primitive
divisor of $P$ which does not contain the whole alternating group,
and hence is not identical with the alternating or the symmetric
group itself, can contain no transposition and no
cyclic permutation of only three figures.

Let $P$ then be the symmetric permutation group of the
$n$ figures $0, 1, 2 \dots n - 1$, and $Q$ a primitive and transitive
divisor of $P$ of index $j$ which does not contain the alternating
group. Further, let $P$ be decomposed into the subsidiary groups:
\weq{4}{P = Q + Q\pi_1 + Q\pi_2 + \dots + Q\pi_{j-1}.}
We consider the system of transpositions:
\weq{5}{(0, 1), (0, 2), \dots (0, n-1),}
none of which can occur in $Q$. Now if $j < n$, then at least
two of these permutations, say $(0,1), (0,2)$, must occur in the
%CONT>>
